\documentclass[12pt]{article}
\usepackage[utf8]{inputenc}
\usepackage[T1]{fontenc}
\usepackage{hyperref}
\title{Mushungulu Language Resources for K-12 Curriculum Development}
\author{Compiled by Claude Sonnet 5.5 (claude-sonnet-5-5) for Andr\'as Kornai}
\date{}
\begin{document}
\maketitle

\section{Introduction}
Mushungulu (ISO 639-3 xma; also Mushunguli) is a Bantu variety spoken in southern Somalia. The search results describe it as a dialect of Zigula (Chizigula), the language of the Zigua of Tanzania. Speaker estimates vary: about 23,000 (Verbix), about 34,000 (Wikipedia) and about 42,000 (Joshua Project). The speakers are thought to descend from people who escaped slavery in north-eastern Tanzania around 1840. Many also speak Somali or Maay. No official orthography was found. One dictionary project notes that some speakers dislike the name ``Mushungulu'' because Somalis often use it disparagingly, so ``Chizigula'' may be preferred. Joshua Project reports that Bible translation has started. I found no school curriculum and no teaching materials.

\section{Primary Educational Resources}
I found no primers, textbooks or school programmes. Materials for Zigula in Tanzania might be adaptable, but I did not find any in this search.

\section{Linguistic and Reference Materials}
\begin{itemize}
\item Somali Chizigula dictionary (Webonary): \url{https://www.webonary.work/somali-chizigula/?p=9444}
\item Zigula language (Wikipedia): \url{https://en.wikipedia.org/wiki/Zigula_language}
\item Zigula language (Wikipedia mirror, code xma): \url{https://wikipedia.com/wiki/ISO_639:xma}
\item Verbix endangered languages entry: \url{https://docs.verbix.com/EndangeredLanguages/Mushungulu}
\item Glottolog reference (Menkhaus 2003): \url{https://glottolog.org/resource/reference/id/27783}
\end{itemize}

\section{Religious and Community Resources}
\begin{itemize}
\item Joshua Project, Zigula in Somalia: \url{https://joshuaproject.net/people_groups/13892/SO}
\item Joshua Project language page for xma: \url{https://joshuaproject.net/index.php/languages/xma}
\item Somali Bantus (Wikipedia): \url{https://www.wikipedia.org/wiki/Mushungulu}
\end{itemize}
I found no community organisations of Mushungulu speakers themselves.

\section{Digital and Audio Resources}
The Webonary dictionary above is the only digital language resource found. I found no audio or video resources.

\section{Curriculum Development Context}
There is no orthography standard, no curriculum and no teaching tradition in Somalia. Development would need community consultation on naming and writing, and most likely work on a shared Zigula standard with Tanzanian materials. The Webonary dictionary is the best starting point.
\end{document}